\section{Introduction}
Choosing a latent dimension is an operational step in neural asset pricing: the researcher must decide how many factors to retain and which pricing errors matter for that choice. The practical question is whether a dimension selected under one evaluation task remains useful under another, and how much pricing fit is sacrificed when it does not. A change in the identity of the minimum is insufficient: a near tie and a large deterioration have different implications.

We distinguish three quantities: concentration in hidden activations, the factor count minimizing a specified empirical pricing loss, and the number of structural priced risks. Flexible predictors and characteristic-based factor models address related but different tasks \citep{kelly2019characteristics,gu2020empirical,gu2021autoencoder,chen2024deep}. We do not claim that these studies mistake prediction for structural identification. Our contribution is a bounded empirical audit of how the first two quantities behave under explicitly stated comparisons.

The original diagnostic networks provide complete candidate loss profiles across three asset spans, six dimensions, five initializations and 21 evaluation weights. We use those frozen profiles to report both preferred dimensions and the cost of carrying the preferred dimension from one endpoint to the other. Cost means additional pricing norm on the same evaluation sample, not a welfare loss or prospective investment outcome. The comparisons are retrospective; the sample minimum is not an oracle population benchmark.

The evidence separates controlled variation from bundled variation. Changing evaluation weights within a fixed model family isolates the sensitivity of the evaluation criterion. The archived characteristic-pipeline comparison also changes construction, coverage and timing, so it cannot identify a feature-deletion effect. A later same-panel deletion control does not reproduce that pipeline reversal. Market hedging and spectral diagnostics provide clues about payoff directions, but wide intervals prevent a precise economic attribution.

A conditional-autoencoder extension tests a boundary of the original narrative. Eight factors, the upper candidate-grid endpoint, win 87.2\% of development resampling draws. The complete archived secondary grid also selects eight in all fifteen autoencoder asset--weight tasks, while the matched linear-loading winners vary. The five-factor candidate has a development pricing norm 11.2\% above the eight-factor candidate; the corresponding excess for linear loadings is only 1.4\%. These are conditional descriptive differences, not significant performance or equivalence findings. A known-truth calibration shows that the new simultaneous intervals undercover, and we do not use them to establish equivalence or structural nonidentification.

Related work in this journal provides useful distinctions. \citet{lo2023qf} study risk-premium prediction and local prediction explanations, whereas our object is the interpretation and cost of candidate factor counts. \citet{xu2024qf} show the importance of a demanding benchmark when assessing machine-learning forecasts. \citet{jiang2025qf} develop a factor-model test; we do not propose a replacement test. \citet{dom2025qf} show that covariance-estimator comparisons depend on portfolio construction. Our related concern is the dependence of latent-dimension comparisons on evaluation design, without claiming that a change of objective producing a different optimum is itself a new theorem.

The study is limited to the archived model families and historical data. The original networks failed earlier benchmark-advancement criteria; the conditional-autoencoder control is exploratory and does not access 2020--2025. Simulations demonstrate possible mechanisms, not their empirical causal prevalence. The following sections define the objects, describe the data and chronology, report candidate losses and controls, and assess the narrower role of representation geometry. The separate updating replay remains in the previously archived supplement and is outside the main research question.
